\subsection{Initial reductions for nineteen-arrow triangles}

\begin{lemma}[Proved: a directed graph as a right-zero action category]
\label{CP2:A-directed-graph-right-zero-realization}
Let $\Gamma$ be a finite nonempty graph with a chosen direction on
each edge.  Let a monoid $M$ act on the right by maps preserving
vertices, edges, their directions and their two endpoint maps.
There is a finite left-$R_3$, right-$M$ biset $S$ and an actual
natural $M$-equivariant homotopy equivalence
\[
 B(R_3,S)\longrightarrow\Gamma.
\]
Here $R_3=\{1,e_0,e_1\}$ with $e_i e_j=e_j$.
Consequently the classifying space of the corresponding two-object
collage is homotopy equivalent to
$\operatorname{hocofib}(\Gamma_{hM}\to BM)$.
\end{lemma}
\begin{proof}
Take $S=\operatorname{Vert}(\Gamma)\sqcup\operatorname{Edge}(\Gamma)$.
Both $e_i$ fix every vertex, and $e_i$ sends an edge to its $i$th
endpoint.  Thus $e_i e_j=e_j$ on $S$, and the prescribed $M$ action
commutes with this left action.

Let $G$ be the graph category with one object for each vertex and
each edge, and two arrows from an edge object to its two endpoints.
The arrows remain distinct when the endpoints coincide.  Its nerve
is the barycentric subdivision of $\Gamma$.
The action category $E=(R_3,S)$ has the same objects; its vertex
endomorphism monoid is $R_3$, its edge endomorphisms are identities,
and its two arrows out of an edge are the two endpoint actions.
Send every vertex endomorphism to an identity and each endpoint
arrow to the corresponding arrow of $G$.  This defines an actual
functor $F:E\to G$, natural under every prescribed graph map.

For an edge object $g$ of $G$, $(F\downarrow g)$ is a point.
For a vertex $v$, the comma category has a central object with
endomorphism monoid $R_3$, and one additional object for each
incident endpoint arrow.  That additional object has a unique
arrow to the central object, fixed by every subsequent central
endomorphism.  The functor to the central copy of $R_3$ sends
this arrow of endpoint type $i$ to $e_i$.  It is a retraction,
and the endpoint arrows together with the central identity give
a natural transformation from the identity to its composite with
the inclusion.  The equation needed for naturality is
$ne_i=e_i$ for $n\in R_3$.
Hence this comma category has the homotopy type of $BR_3$, which
is contractible.  Quillen's Theorem~A proves that $BF$ is a
homotopy equivalence.  The construction is $M$-equivariant, so
the monoid homotopy-orbit bar gives the same equivalence after
taking $M$ homotopy orbits.  Finally $BR_3$ is contractible;
the two-block homotopy-pushout model gives the asserted cofiber.
\end{proof}

\begin{theorem}[Proved: directed graph cofibers through source order eight]
\label{CP2:A-directed-graph-eight}
In the preceding lemma, suppose $|M|\leq8$ and
$Y=\operatorname{hocofib}(\Gamma_{hM}\to BM)$ has bounded rational
homology.  All rational degree-two products of $Y$ vanish.
\end{theorem}
\begin{proof}
If $H^2(BM;\mathbb Q)=0$, every second class of $Y$ is relative
to $BM$, hence lifts from the actual suspension quotient.
Products of these classes are zero.
Otherwise the minimal-ideal homology formula requires at least
two rows and two columns.  A nontrivial minimal group would
give at least eight ideal elements and a distinct global identity,
contradicting $|M|\leq8$.  A group itself has no positive rational
cohomology.  Thus the minimal ideal is a proper rectangular band,
with shape $2\times2$, $2\times3$, or $3\times2$.

Realize $Y$ by the actual $R_3$ collage of the preceding lemma.
For the first two shapes use, respectively,
Theorems~\ref{CP2:A-two-by-two-band-right-zero-all-states} and
\ref{CP2:A-two-by-three-band-general-right-zero}.
For the third shape there are at most two outside elements;
Theorem~\ref{CP2:A-two-column-two-outside-graph} applies to the
bounded actual collage and kills all products.
These statements allow every finite state count, so the number
of vertices and edges in the chosen realization is unrestricted.
\end{proof}

\begin{corollary}[Proved: a semisimple pair opposite a corner through order eight]
\label{CP2:A-semisimple-pair-eight}
Let a three-object category have full linear acyclic support.
If its first two rational corner algebras are semisimple and its
last corner has order at most eight, bounded rational homology
implies that every rational degree-two product is zero.
The opposite statement holds as well.  In particular
$a,b\leq2$, $c\leq8$ excludes every projective-plane profile,
independently of the three cross-hom sizes.
\end{corollary}
\begin{proof}
The actual graph comparison of
Lemma~\ref{CP2:A-baseline-semisimple-pair-graph} has two vertex
classes, one for each source object, and edges directed between
these two classes.  A commuting corner action sends each edge
to an edge and preserves its two endpoint classes.  Thus the
comparison has exactly the direction-preserving hypothesis of
the preceding theorem, after taking opposites if needed.

Write $D$ for the full source pair and collapse $BD$ in the
two-block model.  Since $H_j(BD;\mathbb Q)=0$ for $j\geq2$,
the resulting actual cofiber $Y$ has bounded rational homology.
The map $BC\to Y$ induces a surjection
$H^2(Y;\mathbb Q)\to H^2(BC;\mathbb Q)$.
The natural graph comparison gives the rational ring comparison
with the cofiber of the directed graph action.  Its products
vanish by the theorem, so their pullbacks vanish as well.
Monoids of order one or two have semisimple rational algebras.
\end{proof}

\begin{corollary}[Proved: a small initial singleton at nineteen arrows]
\label{CP2:A-nineteen-small-initial-singleton}
A full-linear three-object category of nineteen arrows cannot
have type $\mathbb CP^2$ if
\[
 a\leq4,\qquad x=1,\qquad c\leq11.
\]
The opposite condition $c\leq4$, $y=1$, $a\leq11$ also excludes it.
\end{corollary}
\begin{proof}
The small-corner dichotomy used in
Corollary~\ref{CP2:A-eighteen-small-initial-singleton} is independent
of the total arrow budget: either the initial corner has
contractible classifying space, or its trivial left rational
module is projective by averaging its nontrivial minimal group.
In the first case Lemma~\ref{CP2:A-linear-contractible-singleton}
deletes the middle object by an actual homotopy equivalence.
It removes its corner and the two adjacent hom-sets, at least
three arrows, contradicting the already proved exclusion through
eighteen arrows.  In the second case
Corollary~\ref{CP2:A-eleven-projective-singleton} kills the products.
The dual argument proves the opposite assertion.
\end{proof}

\begin{proposition}[Proved: the order-eight middle, three-by-five singleton triangle]
\label{CP2:A-nineteen-three-five-singleton-triangle}
The profile $(a,b,c;x,y,z)=(1,8,1;3,5,1)$, and its opposite,
cannot have classifying space of type $\mathbb CP^2$.
\end{proposition}
\begin{proof}
The actual compression in
Lemma~\ref{CP2:A-singleton-triangle-monoid} would give a
twenty-three-element monoid
\[
 P=M\sqcup(V\times W),\qquad |M|=8,\quad |V|=3,\quad |W|=5,
\]
whose minimal ideal is exactly the displayed $3\times5$ band.
Its complement $M$ is a submonoid, by the literal multiplication
in that compression.

The complete minimum-monoid classification
Theorem~\ref{A:monoid-minimum23-classification} leaves exactly
the checked monoid and its opposite.  In the former's literal
table, its minimal band is $\{1,\ldots,15\}$, the identity is $0$,
and its complement is $\{0,16,17,18,19,20,21,22\}$.
The product $16\cdot18=1$ lies in the band, so this complement
is not a submonoid.  The opposite has the same defect with
the two factors interchanged.  The minimal ideal is characteristic
under monoid isomorphisms, so the submonoid-complement property
cannot change under a relabelling.  This contradicts the required
property of $P$ in either orientation.
The complete literal binding, including every outside product
landing in the band, is
\path{converse/scope19/cp23_complement_obstruction.json}.
\end{proof}
